% Printed142; principalPDF149: blank verso.
\par\clearpage\null\thispagestyle{empty}\clearpage
% Printed143; principalPDF150.
\begin{center}
{\Large\textbf{POSTHUMOUS PAPERS.}}\par\medskip
{\large\textbf{[CONFORMAL MAPPING OF THE SPHEROID\\IN THE PLANE.]}}\par\smallskip
\rule{18mm}{.4pt}\par\medskip
[1.]\par\smallskip
[Calculation of geographic latitude and longitude from plane rectangular coordinates.]
\end{center}
[In the conformal mapping of the surface of the spheroid onto the plane, one meridian, the principal meridian, is represented by a straight line, the \(x\)-axis. Every segment of the abscissa axis is equal in length to the portion of the principal meridian of which it is the image.

Let \(a\) be the semimajor axis and \(e\) the eccentricity of the meridian ellipse. Let a point on the surface of the ellipsoid with latitude \(\Phi\) and longitude \(\lambda\) correspond to the rectangular coordinates \(x,y\) in the plane. Let the intersection of the parallel at latitude \(\Phi\) with the principal meridian have abscissa \(\xi\), and let the endpoint of the abscissa \(x\) have latitude \(\varphi\). Thus \(\xi\) and \(x\) are also the corresponding meridian arcs measured from the equator.]

Let a line element be
\[=\sqrt{\{\mathrm d\Phi^2+(\Theta(\Phi))^2\mathrm d\lambda^2\}};\]
\noindent and set
\[\int\frac{\mathrm d\Phi}{\Theta(\Phi)}=\mathrm f(\xi),\]

% Printed144; principalPDF151.
\par\clearpage
\noindent144\hfill{\small\textbf{POSTHUMOUS PAPERS.}}\hfill\null\par\medskip
\noindent then
\[\mathrm f(x+iy)=\mathrm f(\xi)+i\lambda=\mathrm F(\Phi)+i\lambda.\]
[For the spheroid,
\[
\begin{aligned}
\mathrm ds^2&=\frac{aa(1-ee)^2}{(1-ee\sin\Phi^2)^3}\mathrm d\Phi^2+\frac{aa\cos\Phi^2}{1-ee\sin\Phi^2}\mathrm d\lambda^2\\
&=\frac{aa(1-ee)^2}{(1-ee\sin\Phi^2)^3}\GaussOldenburgBracePair{\mathrm d\Phi^2+\left(\frac{(1-ee\sin\Phi^2)\cos\Phi}{1-ee}\right)^2\mathrm d\lambda^2}.
\end{aligned}
\]
Expanding \(\mathrm f(x+iy)\) in powers of \(iy\) gives:
\[
\begin{aligned}
\mathrm f(\xi)=\mathrm F(\Phi)&=\mathrm f(x)-\tfrac12yy\mathrm f''(x)+\tfrac1{24}y^4\mathrm f^{\mathrm{IV}}(x)-\ldots\\
\lambda&=y\mathrm f'(x)-\tfrac16y^3\mathrm f'''(x)+\tfrac1{120}y^5\mathrm f^{\mathrm V}(x)-\ldots.
\end{aligned}
\]
For \(y=0\), we have \(\Phi=\varphi\), so]
\[\mathrm F(\varphi)=\mathrm f(x).\]
[Since \(\dfrac1{\Theta(\varphi)}=\mathrm F'(\varphi)\), we have]
\[\mathrm F'(\varphi)=\frac{1-ee}{(1-ee\sin\varphi^2)\cos\varphi}.\]
[\(x\) increases with \(\varphi\), so]
\[\mathrm dx=\frac{a(1-ee)\mathrm d\varphi}{(1-ee\sin\varphi^2)^{\frac32}},\]
\noindent[hence]
\[\mathrm f'(x)=\frac{(1-ee\sin\varphi^2)^{\frac12}}{a\cos\varphi},\]
\noindent[and further]
\[
\begin{aligned}
\mathrm f''(x)&=\frac{(1-ee\sin\varphi^2)\sin\varphi}{aa\cos\varphi^2}\\
\mathrm f'''(x)&=\frac{(1-ee\sin\varphi^2)^{\frac32}}{a^3(1-ee)\cos\varphi^3}\{1+(1-3ee)\sin\varphi^2+ee\sin\varphi^4\}\\
&\left[=\frac{(1-ee\sin\varphi^2)^{\frac32}}{a^3(1-ee)\cos\varphi^3}\{(1-ee)(2-\cos\varphi^2)+ee\cos\varphi^4\}\right]\\
\mathrm f^{\mathrm{IV}}(x)&=\frac{(1-ee\sin\varphi^2)^2\sin\varphi}{a^4(1-ee)^2\cos\varphi^4}\{5-9ee+(1-4ee+15e^4)\sin\varphi^2\\
&\hspace{32mm}+(ee-13e^4)\sin\varphi^4+4e^4\sin\varphi^6\}\\
&\left[=\frac{(1-ee\sin\varphi^2)^2\sin\varphi}{a^4(1-ee)^2\cos\varphi^4}\{(1-ee)^2(6-\cos\varphi^2)+ee(1-ee)\cos\varphi^4-4e^4\cos\varphi^6\}\right]
\end{aligned}
\]
{\centering etc.\par}
[The rule for differentiation is as follows.]

% Printed145; principalPDF152.
\par\clearpage
\noindent\null\hfill{\small\textbf{CONFORMAL MAPPING OF THE SPHEROID IN THE PLANE.}}\hfill145\par\medskip
Set
\[\frac{\sqrt{(1-ee\sin\varphi^2)}}{a\cos\varphi}=t,\qquad[\sin\varphi=s]\]
\noindent and
\[\mathrm f^n(x)=t^nu,\]
\noindent then
\[\mathrm f^{n+1}(x)=t^{n+1}\left\{nus+\frac{\partial u}{\partial s}.\frac{(1-ss)(1-eess)}{1-ee}\right\}\]
\noindent or, if powers of \(ee\) are neglected,
\[\mathrm f^{n+1}(x)=t^{n+1}\left\{nus+\frac{\partial u}{\partial s}(1-ss)(1+ee(1-ss))\right\}.\]
This gives
\[
\begin{array}{c|c}
n&u\\\hline
1&1\\2&s\\3&1+ss+ee(1-ss)^2\\4&5s+s^3+ee(1-ss)^2s
\end{array}
\]
[Substituting the values of \(\mathrm f'(x)\) and \(\mathrm f'''(x)\) into the equation for \(\lambda\) gives:]
\[
\lambda=\frac{(1-ee\sin\varphi^2)^{\frac12}}{a\cos\varphi}y-\frac{(1-ee\sin\varphi^2)^{\frac32}}{6a^3(1-ee)\cos\varphi^3}(2(1-ee)-(1-ee)\cos\varphi^2+ee\cos\varphi^4)y^3\ldots
\]
[If we set]
\[\mathrm F(\Phi)=\mathrm F(\varphi)+\omega,\]
\noindent[where, consequently,
\[\omega=-\tfrac12yy\mathrm f''(x)+\tfrac1{24}y^4\mathrm f^{\mathrm{IV}}(x)\ldots\]
\noindent then]
\[\begin{aligned}\Phi&=\mathrm G\{\mathrm F(\varphi)+\omega\}\\&=\varphi+A\omega+B\omega\omega+C\omega^3+\ldots,\end{aligned}\]
\noindent[where]
\[A=\left[\frac{\mathrm d\varphi}{\mathrm dF}=\right]\frac{(1-ee\sin\varphi^2)\cos\varphi}{1-ee}.\]
[If we set
\[-\mathrm F'(\varphi)=\frac{\partial\omega}{\partial\varphi}\]
\noindent then, however,]
\[
\begin{aligned}
&1+\frac{\partial A}{\partial\varphi}.\omega+\frac{\partial B}{\partial\varphi}.\omega\omega+\frac{\partial C}{\partial\varphi}.\omega^3\ldots\\
&=A\mathrm F'+2B\mathrm F'.\omega+[3C\mathrm F'.\omega\omega+4D\mathrm F'.\omega^3+]\ldots;
\end{aligned}
\]
\noindent{\small IX.\hfill19}\par

% Printed146; principalPDF153.
\par\clearpage
\noindent146\hfill{\small\textbf{POSTHUMOUS PAPERS.}}\hfill\null\par\medskip
\noindent[it follows that]
\[
\begin{array}{cc}
A\mathrm F'=1&2B=AA'\\2B\mathrm F'=A'&3C=AB'\\3C\mathrm F'=B'&4D=AC'\\\text{etc.}&\text{etc.}
\end{array}
\]
\[
\begin{aligned}
A'&=-\frac{\sin\varphi}{1-ee}(1+2ee-3ee\sin\varphi^2)\\
B&=-\frac{1-ee\sin\varphi^2}{2(1-ee)^2}\sin\varphi\cos\varphi(1+2ee-3ee\sin\varphi^2).
\end{aligned}
\]
[In the equation \(\Phi=\varphi+A\omega+B\omega\omega+\ldots\), substituting the value of \(\omega\) first gives
\[\Phi=\varphi-\tfrac12A\mathrm f''(x).yy+\left(\tfrac1{24}A\mathrm f^{\mathrm{IV}}(x)+\tfrac14B(\mathrm f''(x))^2\right)y^4\ldots;\]
\noindent so with the above values of \(A\) and \(B\) and the values of \(\mathrm f''(x)\) and \(\mathrm f^{\mathrm{IV}}(x)\) on p.~144, we obtain:]
\[
\begin{aligned}
\Phi={}&\varphi-\frac{(1-ee\sin\varphi^2)^2\operatorname{tang}\varphi}{2aa(1-ee)}yy\\
&+\frac{(1-ee\sin\varphi^2)^3\sin\varphi}{24a^4(1-ee)^3\cos\varphi^3}\{5-9ee-(2+7ee-21e^4)\sin\varphi^2\\
&\hspace{45mm}+(10ee-22e^4)\sin\varphi^4+4e^4\sin\varphi^6\}y^4\ldots
\end{aligned}
\]
\noindent[or]
\[
\begin{aligned}
\Phi={}&\varphi-\frac{(1-ee\sin\varphi^2)^2\operatorname{tang}\varphi}{2aa(1-ee)}yy\\
&+\frac{(1-ee\sin\varphi^2)^3\sin\varphi}{24a^4(1-ee)^3\cos\varphi^3}\{3(1-ee)^2+(2-11ee)(1-ee)\cos\varphi^2\\
&\hspace{45mm}+10ee(1-ee)\cos\varphi^4-4e^4\cos\varphi^6\}y^4\ldots
\end{aligned}
\]
[If we set \(\Phi=\varphi-z\), then
\[\log\operatorname{hyp}\sin\Phi=\log\operatorname{hyp}\sin\varphi-\frac1{\operatorname{tang}\varphi}z-\tfrac12\frac1{\sin\varphi^2}zz\ldots;\]
\noindent hence:]
\[
\begin{aligned}
\log\operatorname{hyp}\sin\Phi={}&\log\operatorname{hyp}\sin\varphi-\frac{(1-ee\sin\varphi^2)^2}{2aa(1-ee)}yy\\
&+\frac{(1-ee\sin\varphi^2)^3}{12a^4(1-ee)^2}\left\{1-7ee+5ee\cos\varphi^2-\frac{2e^4}{1-ee}\cos\varphi^4\right\}y^4\ldots.
\end{aligned}
\]
{\centering [2.]\par
[Calculation of meridian convergence from plane rectangular coordinates.]\par}
[The meridian convergence \(c\) is the angle that the image of the meridian makes with the line parallel to the \(x\)-axis. Let \(x\) be positive northward. Then:
\[\operatorname{tang}c=-\frac{\mathrm dy}{\mathrm dx}.\]

% Printed147; principalPDF154.
\par\clearpage
\noindent\null\hfill{\small\textbf{CONFORMAL MAPPING OF THE SPHEROID IN THE PLANE.}}\hfill147\par\medskip
From the equation that represents the meridian of longitude \(\lambda\) in the plane,
\[\lambda=\operatorname{const.}=y\mathrm f'(x)-\tfrac16y^3\mathrm f'''(x)\ldots,\]
\noindent it follows, however, that
\[\frac{\mathrm dy}{\mathrm dx}=-\frac{y\mathrm f''(x)-\tfrac16y^3\mathrm f^{\mathrm{IV}}(x)\ldots}{\mathrm f'(x)-\tfrac12yy\mathrm f'''(x)\ldots},\]
\noindent so that:]
\[
\begin{aligned}
\operatorname{tang}c&=\frac{y\mathrm f''(x)-\tfrac16y^3\mathrm f^{\mathrm{IV}}(x)\ldots}{\mathrm f'(x)-\tfrac12yy\mathrm f'''(x)\ldots}\\
&\left[=y\frac{\mathrm f''(x)}{\mathrm f'(x)}\left\{1-yy\left(\tfrac16\frac{\mathrm f^{\mathrm{IV}}(x)}{\mathrm f''(x)}-\tfrac12\frac{\mathrm f'''(x)}{\mathrm f'(x)}\right)\ldots\right\}. \right.
\end{aligned}
\]
Hence, by p.~144, we have]
\[\operatorname{tang}c=D.\frac{\sqrt{(1-ee\sin\varphi^2)}}a y\operatorname{tang}\varphi,\]
\noindent where
\[
\begin{aligned}
\log\operatorname{hyp}D&=\left(-\tfrac16\frac{\mathrm f^{\mathrm{IV}}(x)}{\mathrm f''(x)}+\tfrac12\frac{\mathrm f'''(x)}{\mathrm f'(x)}\right)yy\ldots\\
&=-\tfrac13yy\frac{(1-ee\sin\varphi^2)^2}{aa(1-ee)^2}(1-3ee+2ee\sin\varphi^2)\ldots
\end{aligned}
\]
\noindent[or approximately]
\[=-\tfrac13\frac{yy}{aa}(1-ee).\]
[We also have:]
\[c=y\frac{\mathrm f''(x)}{\mathrm f'(x)}-\tfrac16y^3\left\{\frac{\mathrm f^{\mathrm{IV}}(x)}{\mathrm f'(x)}-\frac{3\mathrm f''(x)\mathrm f'''(x)}{(\mathrm f'(x))^2}+2\left(\frac{\mathrm f''(x)}{\mathrm f'(x)}\right)^3\right\}\ldots\]
\noindent[or, if we set
\[\frac{\mathrm f''(x)}{\mathrm f'(x)}=\mathrm h(x)\]
\noindent then
\[c=y\mathrm h(x)-\tfrac16y^3\mathrm h''(x)\ldots].\]
\[
\begin{aligned}
\mathrm h(x)&=\frac{(1-ee\sin\varphi^2)^{\frac12}\sin\varphi}{a\cos\varphi}\\
\mathrm h'(x)&=\frac{1-ee\sin\varphi^2}{aa(1-ee)\cos\varphi^2}\{1-2ee\sin\varphi^2+ee\sin\varphi^4\}\\
{[}\mathrm h''(x)&=\frac{2(1-ee\sin\varphi^2)^{\frac32}\sin\varphi}{a^3(1-ee)^2\cos\varphi^3}\{1-3ee+(2ee+4e^4)\sin\varphi^2\\
&\hspace{37mm}-(ee+5e^4)\sin\varphi^4+2e^4\sin\varphi^6\}\\
&=\frac{2(1-ee\sin\varphi^2)^{\frac32}\sin\varphi}{a^3(1-ee)^2\cos\varphi^3}\{(1-ee)^2-ee(1-ee)\cos\varphi^4-2e^4\cos\varphi^6\}].
\end{aligned}
\]
\noindent\hfill{\small19*}\par

% Printed148; principalPDF155.
\par\clearpage
\noindent148\hfill{\small\textbf{POSTHUMOUS PAPERS.}}\hfill\null\par\medskip
Thus:
\[
\begin{aligned}
c={}&\frac{\sqrt{(1-ee\sin\varphi^2)}}a\operatorname{tang}\varphi.y\\
&-\tfrac13\frac{(1-ee\sin\varphi^2)^{\frac32}\sin\varphi}{a^3(1-ee)^2\cos\varphi^3}\{1-3ee+(2ee+4e^4)\sin\varphi^2\\
&\hspace{44mm}-(ee+5e^4)\sin\varphi^4+2e^4\sin\varphi^6\}y^3\ldots
\end{aligned}
\]
\noindent[or]
\[
\begin{aligned}
c={}&\frac{\sqrt{(1-ee\sin\varphi^2)}}a\operatorname{tang}\varphi.y\\
&-\tfrac13\frac{(1-ee\sin\varphi^2)^{\frac32}\sin\varphi}{a^3(1-ee)^2\cos\varphi^3}\{(1-ee)^2\quad *\quad-ee(1-ee)\cos\varphi^4-2e^4\cos\varphi^6\}y^3\ldots.
\end{aligned}
\]
[Introducing \(\lambda\), we also have:]
\[c=\lambda\sin\varphi-\frac{(1-ee\sin\varphi^2)^{\frac32}\operatorname{tang}\varphi}{6a^3(1-ee)^2}\{(1-ee)^2-3ee(1-ee)\cos\varphi^2-4e^4\cos\varphi^4\}y^3\ldots.\]
\medskip
{\centering [3.]\par}
The following formulas are the most convenient for numerical calculation. Let
\[
\begin{aligned}
L&=\frac{(1-ee\sin\varphi^2)^{\frac12}y}{a\rho\cos\varphi}\\
M&=\frac{(1-ee\sin\varphi^2)^2\operatorname{tang}\varphi.yy}{2aa(1-ee)\rho}\\
N&=\frac{(1-ee\sin\varphi^2)^{\frac12}}{a\rho}\operatorname{tang}\varphi.y.
\end{aligned}
\]
Then
\[
\begin{aligned}
\log\lambda&=\log L-A.LL\\
\log(\varphi-\Phi)&=\log M-B.LL\\
\log c&=\log N-C.LL,
\end{aligned}
\]
\noindent where the logarithms of \(A,B,C\) are most conveniently taken from a separate table with argument \(\varphi\).

The following formulas are used to calculate this table:
\[
\begin{aligned}
A&=H\left\{0{,}75-\tfrac14\cos2\varphi+\frac{ee}{2(1-ee)}\cos\varphi^4\right\}\\
B&=H\left\{0{,}75+\frac{2-11ee}{4(1-ee)}\cos\varphi^2+\frac{5ee}{2(1-ee)}\cos\varphi^4-\frac{e^4}{(1-ee)^2}\cos\varphi^6\right\}\\
C&=H\left\{1-\frac{ee}{1-ee}\cos\varphi^4-\frac{2e^4}{(1-ee)^2}\cos\varphi^6\right\},
\end{aligned}
\]

% Printed149; principalPDF156.
\par\clearpage
\noindent\null\hfill{\small\textbf{CONFORMAL MAPPING OF THE SPHEROID IN THE PLANE.}}\hfill149\par\medskip
\noindent where
\[H=\tfrac13k\rho\rho.10^7,\qquad\log H=5{,}531\,8128-10\]
[\(k=\) modulus of the Briggsian logarithms and \(\rho=\dfrac1{206264{,}8\ldots}\)

The corrections \(A.LL,B.LL,C.LL\) are obtained in units of the seventh decimal place. \(\lambda,\varphi-\Phi\) and \(c\) are obtained in arcseconds.

The following table of the logarithms of \(A,B,C\) is based on the flattening value \(\frac1{302{,}68}\)].
\begin{center}
\begingroup\scriptsize\setlength{\tabcolsep}{2pt}\renewcommand{\arraystretch}{1.15}
\begin{tabular}[t]{|r|r|r|r|}\hline
\(\varphi\)&\(\log A\)&\(\log B\)&\(\log C\)\\\hline\hline
&\(-10\)&\(-10\)&\(-10\)\\
\(50^\circ\;0'\)&\(5{,}431\,6218\)&\(5{,}510\,4149\)&\(5{,}531\,3175\)\\
\(3\)&\(8555\)&\(2206\)&\(95\)\\
\(6\)&\(5{,}432\,0889\)&\(0263\)&\(5{,}531\,3216\)\\
\(9\)&\(3222\)&\(5{,}509\,8319\)&\(36\)\\
\(12\)&\(5552\)&\(6376\)&\(57\)\\
\(15\)&\(7881\)&\(4432\)&\(77\)\\
\hline
\(18\)&\(5{,}433\,0208\)&\(2488\)&\(97\)\\
\(21\)&\(2532\)&\(0544\)&\(5{,}531\,3318\)\\
\(24\)&\(4854\)&\(5{,}508\,8600\)&\(38\)\\
\(27\)&\(7174\)&\(6656\)&\(59\)\\
\(30\)&\(9494\)&\(4712\)&\(79\)\\
\hline
\(33\)&\(5{,}434\,1810\)&\(2768\)&\(99\)\\
\(36\)&\(4124\)&\(0823\)&\(5{,}531\,3419\)\\
\(39\)&\(6438\)&\(5{,}507\,8878\)&\(39\)\\
\(42\)&\(8748\)&\(6934\)&\(59\)\\
\(45\)&\(5{,}435\,1057\)&\(4989\)&\(79\)\\
\hline
\(48\)&\(3364\)&\(3044\)&\(99\)\\
\(51\)&\(5668\)&\(1100\)&\(5{,}531\,3519\)\\
\(54\)&\(7971\)&\(5{,}506\,9155\)&\(38\)\\
\(57\)&\(5{,}436\,0271\)&\(7210\)&\(58\)\\
\(51\quad0\)&\(2570\)&\(5265\)&\(78\)\\
\hline
\(3\)&\(4867\)&\(3320\)&\(98\)\\
\(6\)&\(7161\)&\(1375\)&\(5{,}531\,3617\)\\
\(9\)&\(9454\)&\(5{,}505\,9430\)&\(37\)\\
\(12\)&\(5{,}437\,1745\)&\(7485\)&\(56\)\\
\(15\)&\(4033\)&\(5540\)&\(76\)\\
\hline
\(18\)&\(6319\)&\(3595\)&\(95\)\\
\(21\)&\(8604\)&\(1650\)&\(5{,}531\,3714\)\\
\(24\)&\(5{,}438\,0886\)&\(5{,}504\,9704\)&\(34\)\\
\(27\)&\(3166\)&\(7759\)&\(53\)\\
\(30\)&\(5444\)&\(5814\)&\(72\)\\
\hline
\(33\)&\(7720\)&\(3869\)&\(91\)\\
\(36\)&\(9994\)&\(1924\)&\(5{,}531\,3810\)\\
\(39\)&\(5{,}439\,2266\)&\(5{,}503\,9979\)&\(30\)\\
\(42\)&\(4535\)&\(8034\)&\(49\)\\
\(45\)&\(6803\)&\(6089\)&\(68\)\\
\hline\end{tabular}
\hfill
\begin{tabular}[t]{|r|r|r|r|}\hline
\(\varphi\)&\(\log A\)&\(\log B\)&\(\log C\)\\\hline\hline
&\(-10\)&\(-10\)&\(-10\)\\
\(51^\circ45'\)&\(5{,}439\,6803\)&\(5{,}503\,6089\)&\(5{,}531\,3868\)\\
\(48\)&\(9068\)&\(4144\)&\(87\)\\
\(51\)&\(5{,}440\,1332\)&\(2199\)&\(5{,}531\,3906\)\\
\(54\)&\(3593\)&\(0254\)&\(24\)\\
\(57\)&\(5852\)&\(5{,}502\,8310\)&\(43\)\\
\(52\quad0\)&\(8109\)&\(6365\)&\(62\)\\
\hline
\(3\)&\(5{,}441\,0364\)&\(4420\)&\(81\)\\
\(6\)&\(2616\)&\(2476\)&\(99\)\\
\(9\)&\(4867\)&\(0531\)&\(5{,}531\,4018\)\\
\(12\)&\(7115\)&\(5{,}501\,8587\)&\(36\)\\
\(15\)&\(9361\)&\(6643\)&\(55\)\\
\hline
\(18\)&\(5{,}442\,1605\)&\(4699\)&\(73\)\\
\(21\)&\(3847\)&\(2755\)&\(91\)\\
\(24\)&\(6087\)&\(0812\)&\(5{,}531\,4110\)\\
\(27\)&\(8325\)&\(5{,}500\,8868\)&\(28\)\\
\(30\)&\(5{,}443\,0560\)&\(6925\)&\(46\)\\
\hline
\(33\)&\(2793\)&\(4982\)&\(64\)\\
\(36\)&\(5024\)&\(3039\)&\(82\)\\
\(39\)&\(7254\)&\(1096\)&\(5{,}531\,4200\)\\
\(42\)&\(9480\)&\(5{,}499\,9153\)&\(18\)\\
\(45\)&\(5{,}444\,1705\)&\(7211\)&\(36\)\\
\hline
\(48\)&\(3927\)&\(5269\)&\(54\)\\
\(51\)&\(6148\)&\(3327\)&\(72\)\\
\(54\)&\(8366\)&\(1384\)&\(89\)\\
\(57\)&\(5{,}445\,0581\)&\(5{,}498\,9443\)&\(5{,}531\,4307\)\\
\(53\quad0\)&\(2795\)&\(7501\)&\(25\)\\
\hline
\(3\)&\(5006\)&\(5560\)&\(43\)\\
\(6\)&\(7215\)&\(3618\)&\(60\)\\
\(9\)&\(9422\)&\(1677\)&\(78\)\\
\(12\)&\(5{,}446\,1627\)&\(5{,}497\,9737\)&\(95\)\\
\(15\)&\(3829\)&\(7796\)&\(5{,}531\,4413\)\\
\hline
\(18\)&\(6029\)&\(5856\)&\(30\)\\
\(21\)&\(8227\)&\(3916\)&\(48\)\\
\(24\)&\(5{,}447\,0423\)&\(1976\)&\(65\)\\
\(27\)&\(2617\)&\(0036\)&\(83\)\\
\(30\)&\(4808\)&\(5{,}496\,8097\)&\(5{,}531\,4500\)\\
\hline\end{tabular}
\hfill
\begin{tabular}[t]{|r|r|r|r|}\hline
\(\varphi\)&\(\log A\)&\(\log B\)&\(\log C\)\\\hline\hline
&\(-10\)&\(-10\)&\(-10\)\\
\(53^\circ30'\)&\(5{,}447\,4808\)&\(5{,}496\,8097\)&\(5{,}531\,4500\)\\
\(33\)&\(6997\)&\(6158\)&\(17\)\\
\(36\)&\(9184\)&\(4219\)&\(34\)\\
\(39\)&\(5{,}448\,1369\)&\(2281\)&\(51\)\\
\(42\)&\(3551\)&\(0343\)&\(68\)\\
\(45\)&\(5731\)&\(5{,}495\,8405\)&\(85\)\\
\hline
\(48\)&\(7909\)&\(6468\)&\(5{,}531\,4602\)\\
\(51\)&\(5{,}449\,0084\)&\(4530\)&\(19\)\\
\(54\)&\(2258\)&\(2594\)&\(35\)\\
\(57\)&\(4428\)&\(0657\)&\(52\)\\
\(54\quad0\)&\(6597\)&\(5{,}494\,8721\)&\(5{,}531\,4669\)\\
\hline
\(3\)&\(8763\)&\(6785\)&\(86\)\\
\(6\)&\(5{,}450\,0927\)&\(4850\)&\(5{,}531\,4702\)\\
\(9\)&\(3089\)&\(2915\)&\(19\)\\
\(12\)&\(5248\)&\(0980\)&\(35\)\\
\(15\)&\(7405\)&\(5{,}493\,9046\)&\(52\)\\
\hline
\(18\)&\(9560\)&\(7112\)&\(68\)\\
\(21\)&\(5{,}451\,1712\)&\(5179\)&\(84\)\\
\(24\)&\(3863\)&\(3245\)&\(5{,}531\,4801\)\\
\(27\)&\(6010\)&\(1313\)&\(17\)\\
\(30\)&\(8156\)&\(5{,}492\,9380\)&\(33\)\\
\hline
\(33\)&\(5{,}452\,0299\)&\(7448\)&\(49\)\\
\(36\)&\(2440\)&\(5516\)&\(65\)\\
\(39\)&\(4579\)&\(3585\)&\(81\)\\
\(42\)&\(6715\)&\(1654\)&\(97\)\\
\(45\)&\(8849\)&\(5{,}491\,9724\)&\(5{,}531\,4913\)\\
\hline
\(48\)&\(5{,}453\,0981\)&\(7794\)&\(29\)\\
\(51\)&\(3110\)&\(5865\)&\(45\)\\
\(54\)&\(5237\)&\(3936\)&\(60\)\\
\(57\)&\(7362\)&\(2007\)&\(76\)\\
\(55\quad0\)&\(9484\)&\(0079\)&\(92\)\\
\hline\end{tabular}
\endgroup
\end{center}

% Printed150; principalPDF157.
\par\clearpage
\noindent150\hfill{\small\textbf{POSTHUMOUS PAPERS.}}\hfill\null\par\medskip
{\centering [4.]\par}
[We also have]
\[\lambda=A'\frac{\sqrt{(1-ee\sin\varphi^2)}}a.\frac y{\cos\varphi},\]
\noindent where
\[\log\operatorname{hyp}A'=-\tfrac13cc-\tfrac16\frac{(1-ee\sin\varphi^2)^2}{aa(1-ee)}yy\ldots\]
\noindent or
\[=-\tfrac13\lambda\lambda+\frac{1-ee}{6aa}\left(1-\frac{e^4}{(1-ee)^2}\cos\varphi^4\right)yy\ldots\]
\noindent[and approximately]
\[=-\tfrac13\lambda\lambda+\frac{1-ee}{6aa}yy.\]
[Furthermore,]
\[\Phi=\varphi-B'\frac{(1-ee\sin\varphi^2)^2}{2aa(1-ee)}\operatorname{tang}\varphi.yy,\]
\noindent where [approximately]
\[
\begin{aligned}
\log\operatorname{hyp}B'&=-\tfrac14\lambda\lambda-\tfrac16\frac{yy}{aa}\\
&=-\tfrac5{12}\lambda\lambda+\tfrac16cc\\
&=-\tfrac14cc-\tfrac5{12}\frac{yy}{aa}.
\end{aligned}
\]
We can also set
\[c=C'\frac{\sqrt{(1-ee\sin\varphi^2)}}a\operatorname{tang}\varphi.y,\]
\noindent where
\[
\log\operatorname{hyp}C'=-\tfrac13cc-\tfrac13\left\{\frac{1-ee\sin\varphi^2}{aa}-\frac{ee(1-ee\sin\varphi^2)\cos\varphi^2(1-ee+2ee\cos\varphi^2)}{aa(1-ee)^2}\right\}yy\ldots
\]
\noindent[or approximately
\[=-\tfrac13cc-\tfrac13\frac{1-ee}{aa}yy].\]

% Printed151; principalPDF158.
\par\clearpage
\noindent\null\hfill{\small\textbf{CONFORMAL MAPPING OF THE SPHEROID IN THE PLANE.}}\hfill151\par\medskip
{\centering [5.]\par}
If the positions of points are represented by coordinates \(x,y\), the former increasing uniformly along the meridian, the latitude \(\Phi\) of a place, the convergence \(c\) of its meridian, and its longitude \(\lambda\) are found as follows.

Let \(\varphi\) be the latitude corresponding to the abscissa \(x\) on the meridian,
\[
\begin{aligned}
q&=\operatorname{Compl.}\log\sqrt{(1-ee\sin\varphi^2)}\\
E&=\frac{1-ee}{6aa}.k.10^7\\
H&=\frac{k.10^7}{3.206265^2}.
\end{aligned}
\]
[\(k\) modulus of the Briggsian logarithms
\[
\begin{aligned}
\log B&=\log\frac{206264{,}8\ldotp\ldotp}a-q,\qquad\log A=\log\frac{206264{,}8\ldotp\ldotp}{a(1-ee)}-3q\\
\log C&=\log\frac{206264{,}8\ldotp\ldotp}{2aa(1-ee)}-4q.
\end{aligned}
\]
For the flattening value \(\frac1{302{,}68}\), we have:]
\[
\begin{aligned}
\log E&=2{,}24756\,29768-10,\qquad\log H=5{,}53181\,27903-10\\
\log B&=8{,}50982\,72984-10-q\\
\log C&=\log\frac{AB}{2.206264{,}8\ldotp\ldotp}=1{,}40707\,38825-10-4q.
\end{aligned}
\]
[To calculate \(A\) and \(B\), the table on p.~84 can be used directly, or \(q\) can be taken from the table on p.~77.]
\[
\begin{aligned}
c&=(1)By\operatorname{tang}\varphi&\log(1)&=-2Eyy-Hcc\\
\lambda&=(2)B\frac y{\cos\varphi}&\log(2)&=+Eyy-H\lambda\lambda\\
\Phi&=\varphi-(3)Cyy\operatorname{tang}\varphi&\log(3)&=-\tfrac34H\lambda\lambda-Eyy\\
&&&=-\tfrac54H\lambda\lambda+\tfrac12Hcc\\
&&&=-\tfrac34Hcc-\tfrac52Eyy.
\end{aligned}
\]
[\(\log(1),\log(2),\log(3)\) are obtained in units of the seventh decimal place. \(c,\lambda,\Phi-\varphi\) are obtained in arcseconds.]
